% =====================================================================
%  MLE-AI-201 Machine Learning -- Final Project report
%  IEEE-style two-column layout using standard LaTeX packages only.
%
%  Build:   tectonic report.tex      (from report/)
%
%  Numbers come from the files written by `python -m src.run_all`; each is
%  followed by a "% src:" comment naming the file and the row/column (or key)
%  it was taken from. Claims about a method are followed by a
%  "% src: file.py:function" comment naming the code that implements it.
% =====================================================================
\documentclass[10pt,twocolumn,a4paper]{article}

% ---- Geometry: IEEE-like margins / two narrow columns ----------------
\usepackage[a4paper,top=0.75in,bottom=1.0in,left=0.68in,right=0.68in,columnsep=0.25in]{geometry}

% ---- Times-like font ---------------------------------------------------
\usepackage{mathptmx}
\usepackage[T1]{fontenc}
\usepackage[utf8]{inputenc}

% ---- Packages --------------------------------------------------------
\usepackage{graphicx}
\usepackage{amsmath,amssymb}
\usepackage{booktabs}
\usepackage{array}
\usepackage{titlesec}
\usepackage{caption}
\usepackage{cite}
\usepackage[hidelinks]{hyperref}
\usepackage{url}

% ---- IEEE-style section headings -------------------------------------
\renewcommand{\thesection}{\Roman{section}}
\renewcommand{\thesubsection}{\Alph{subsection}}
\titleformat{\section}[block]{\centering\normalsize\scshape}{\thesection.}{0.5em}{}
\titleformat{\subsection}[block]{\normalsize\itshape}{\thesubsection.}{0.5em}{}
\titlespacing*{\section}{0pt}{1.4ex plus 1ex minus .2ex}{1ex}
\titlespacing*{\subsection}{0pt}{1.0ex plus 1ex minus .2ex}{0.6ex}

% ---- Abstract / Index-Terms in the IEEE boldface-dash style ----------
\newcommand{\abstractkw}[1]{\noindent\textbf{\textit{#1}}}

\captionsetup{font=small,labelsep=period}
\setlength{\parindent}{1em}

% =====================================================================
\begin{document}

% ------------------------- TITLE BLOCK --------------------------------
\twocolumn[
  \begin{@twocolumnfalse}
  \begin{center}
    {\LARGE\bfseries Pricing Baku Real Estate from Scratch:\\Decision Trees and Pegasos SVMs on bina.az Listings\par}
    \vspace{1.2em}
    % Authors: Surname, Given name --- alphabetical by surname.
    {\large
      Aghazade, Ibrahim\quad Aliyev, Sahib\quad Mammadli, Mahammad\quad Sakhrayev, Khizri\par}
    \vspace{0.5em}
    {\itshape MLE-AI-201 Machine Learning --- AI Academy, National Artificial Intelligence Center\par}
    \vspace{1.4em}
  \end{center}
  \begin{quote}
    \small
    \abstractkw{Abstract---}%
    We predict the prices of real-estate listings scraped from bina.az, using
    91{,}314 cleaned sale listings. % src: results/report_facts.csv [rows_clean]
    Task~A regresses the logarithm of the price; Task~B classifies each
    listing as premium or standard, split at the median training price of
    215{,}000~AZN. % src: results/report_facts.csv [tier_threshold_train_median]
    We implemented two models from scratch in NumPy: a CART-style decision
    tree with Gini, entropy and MSE criteria, and a Pegasos sub-gradient
    support-vector machine, either linear or on random Fourier features that
    approximate an RBF kernel. Both are benchmarked against matched
    scikit-learn models on identical splits, with 5-fold cross-validation, a
    held-out test set and timing. On the test set our regression tree
    reaches $R^2 = 0.853$ on log price, % src: results/benchmark_test_task_a.csv [Tree (ours), R2 log 0.8531]
    against 0.852 for scikit-learn. % src: results/benchmark_test_task_a.csv [Tree (sklearn), R2 log 0.8522]
    Our Gini tree reaches F1 0.926 and ROC-AUC 0.971, % src: results/benchmark_test_task_b.csv [Tree gini (ours), F1 0.9258, ROC-AUC 0.9705]
    and our linear Pegasos SVM, chosen on validation over the Fourier-feature
    variant, reaches F1 0.904. % src: results/benchmark_test_task_b.csv [Pegasos linear (ours), F1 0.9040]
    Each of our models is within 0.007 F1 of its best matched scikit-learn % src: results/error_analysis.md [Findings, bullet 1: within 0.007, largest gap -0.0065]
    baseline, and our trees are nearly identical to scikit-learn's, but they
    fit about 277 times more slowly. % src: 86.7520 / 0.3128 from results/benchmark_test_task_b.csv [Fit s, Tree gini (ours) / (sklearn)]
    Errors concentrate near the tier threshold, at both price extremes and
    in districts with little training data. The main takeaway: careful
    from-scratch implementations reach library accuracy on this data; what
    they lose is speed. Code:
    \href{https://github.com/mmahammad57/binaaz-ml-project}{\texttt{github.com/mmahammad57/binaaz-ml-project}}
    (tag \texttt{v1.0-final}).

    \vspace{0.6em}
    \abstractkw{Index Terms---}decision tree, support-vector machine,
    Pegasos, regression, classification, real-estate price prediction.
  \end{quote}
  \vspace{1.0em}
  \end{@twocolumnfalse}
]

% =====================================================================
\section{Introduction}
Listing prices on bina.az, a property portal for Azerbaijan, are set by
sellers and agents. A model
that predicts a listing's price from its size, type and location helps
both sides judge whether an asking price is reasonable. We study two tasks
on the same listings. \emph{Task~A} is regression of $\log(\text{price})$.
\emph{Task~B} is classification of a derived \emph{price tier}: premium
if the price is at least the median training price, standard otherwise.

Our contributions are:
\begin{itemize}
  \item a decision tree written from scratch in NumPy, for regression (MSE)
  and classification (Gini or entropy), with the usual stopping rules and
  impurity-based feature importances;
  \item a Pegasos support-vector machine written from scratch, linear or on
  random Fourier features (RFF) for a non-linear decision boundary;
  \item a benchmark of both against matched scikit-learn
  models~\cite{sklearn} on identical data: cross-validation, a held-out
  test set, timing and a scaling study;
  \item an error analysis by price range, district, category and distance
  from the tier threshold;
  \item a reproducible pipeline: one command regenerates every number,
  table and figure in this report.
\end{itemize}

% =====================================================================
\section{Background}
\label{sec:background}
\subsection{Decision trees}
A CART tree~\cite{cart} splits a node $S$ on one feature $j$ and a
threshold $t$ into $S_L = \{x \in S : x_j \le t\}$ and
$S_R = S \setminus S_L$, choosing the split that most reduces an impurity
$I$. For classes $k$ with proportions $p_k$ in $S$, and for a regression
target $y$ with mean $\bar{y}_S$,
\begin{gather}
  I_{\text{Gini}}(S) = 1 - \textstyle\sum_k p_k^2, \qquad
  I_{\text{ent}}(S) = -\textstyle\sum_k p_k \log_2 p_k, \label{eq:imp}\\
  I_{\text{MSE}}(S) = \tfrac{1}{|S|}\textstyle\sum_{i \in S} (y_i - \bar{y}_S)^2 .\notag
\end{gather}
The impurity decrease, called information gain for entropy, is
\begin{equation}
  \Delta(S, j, t) = I(S) - \tfrac{|S_L|}{|S|} I(S_L) - \tfrac{|S_R|}{|S|} I(S_R).
  \label{eq:gain}
\end{equation}
Splitting is repeated on each child until a stopping rule holds; a leaf
predicts the class proportions or the mean target of its rows.

\subsection{Support-vector machines and Pegasos}
The soft-margin SVM~\cite{cortes} with labels $y_i \in \{-1, +1\}$
minimises $\tfrac{1}{2}\lVert\mathbf{w}\rVert^2 + C\sum_i \xi_i$ subject
to $y_i(\mathbf{w}^\top\mathbf{x}_i + b) \ge 1 - \xi_i$, $\xi_i \ge 0$.
Eliminating the slacks gives the hinge-loss primal
\begin{equation}
  J(\mathbf{w},b)=\tfrac{\lambda}{2}\lVert\mathbf{w}\rVert^{2}
  +\tfrac{1}{n}\sum_{i=1}^{n}\max\!\big(0,\,1-y_i(\mathbf{w}^{\top}\mathbf{x}_i+b)\big),
  \label{eq:svm}
\end{equation}
which has the same minimiser as the constrained form when
$\lambda = 1/(nC)$. Pegasos~\cite{pegasos} minimises~(\ref{eq:svm}) by
stochastic sub-gradient steps. At step $t$ it draws an example $i$, sets
$\eta_t = 1/(\lambda t)$ and updates
\begin{equation}
  \mathbf{w} \leftarrow (1 - \eta_t\lambda)\,\mathbf{w}
  + \eta_t\,\mathbf{1}\!\left[y_i\mathbf{w}^\top\mathbf{x}_i < 1\right] y_i\mathbf{x}_i ,
  \label{eq:pegasos}
\end{equation}
optionally followed by a projection of $\mathbf{w}$ onto the ball of
radius $1/\sqrt{\lambda}$. After $T$ steps its error in the objective
shrinks like $1/(\lambda T)$, so small $\lambda$ needs many steps.

\subsection{Random Fourier features}
A kernel SVM replaces $\mathbf{x}_i^\top\mathbf{x}_j$ by a kernel such as
the RBF $k(\mathbf{x},\mathbf{x}') = \exp(-\gamma\lVert\mathbf{x}-\mathbf{x}'\rVert^2)$,
whose $n \times n$ Gram matrix is too large for our data.
Rahimi and Recht~\cite{rahimi} approximate it with an explicit map
$\mathbf{z}(\mathbf{x}) = \sqrt{2/D}\,\cos(W^\top\mathbf{x} + \mathbf{a})$,
with $W_{jk} \sim \mathcal{N}(0, 2\gamma)$ and
$a_k \sim \mathcal{U}[0, 2\pi)$, so that
$\mathbb{E}[\mathbf{z}(\mathbf{x})^\top\mathbf{z}(\mathbf{x}')] = k(\mathbf{x},\mathbf{x}')$.
A linear SVM on $\mathbf{z}(\mathbf{x})$ then approximates the RBF-kernel SVM.

% =====================================================================
\section{Data and Preprocessing}
\label{sec:data}
The data is the Kaggle bina.az sale collection~\cite{kaggle}: one CSV of
100{,}775 listings % src: results/report_facts.csv [raw_rows]
and 51 columns. % src: results/report_facts.csv [raw_columns]
All preprocessing is in \texttt{data\_prep.py} (Table~\ref{tab:data}).

{\raggedright\textit{Columns removed.} The pipeline drops 30 columns % src: results/report_facts.csv [dropped_columns]
listed in \texttt{LEAKAGE\_COLUMNS}. % src: data_prep.py:LEAKAGE_COLUMNS, data_prep.py:clean (step 1)
Only two are derived from the target: \texttt{unit\_price}, the price per
m\textsuperscript{2}, and \texttt{total\_price}, which repeats the price;
either would give the answer away. The others carry no information about
the property or would identify listings and sellers: identifiers and URLs
(\texttt{id\_x}, \texttt{id\_y}, \texttt{estate\_id},
\texttt{estate\_details\_id\_x/\_y}, \texttt{rel\_url},
\texttt{estate\_rel\_url}, \texttt{estate\_rel\_url\_x/\_y},
\texttt{img\_url}); scrape and update times (\texttt{datetime\_scrape\_x/\_y},
\texttt{day\_x/\_y}, \texttt{hour\_x/\_y}, \texttt{updated},
\texttt{city\_when}); seller names (\texttt{owner\_name},
\texttt{owner\_title}, \texttt{shop\_name}, \texttt{shop\_title}); free
text (\texttt{address}, \texttt{description}, \texttt{attributes},
\texttt{extra\_info}); and the currency columns
(\texttt{currency\_x/\_y}).\par}

\textit{Columns kept.} The remaining 21 columns are: % src: results/report_facts.csv [raw_columns_kept]
the price (target); the location name (district, metro-station area or
settlement) and city; floor area (m\textsuperscript{2}), number of rooms,
floor (``floor / storeys''), land area (in sot); latitude, longitude and
page views; category (new build, old build, house or villa, office),
building type, renovation, title deed and mortgage availability (all
originally in Azerbaijani); \texttt{products\_label} (agency or
residential complex); and five flag columns that hold a single value or
are missing: \texttt{repair} (renovated), \texttt{bill\_of\_sale} (title
deed available), \texttt{mortgage} (mortgage available), \texttt{vip} and
\texttt{featured} (paid promotion of the listing).

\textit{Cleaning}, in the order of \texttt{clean()}: % src: data_prep.py:clean
(1)~the floor string ``$k$ / $m$'' is split at ``/'' into \texttt{floor}
$=k$ and \texttt{total\_floors} $=m$, unparsable or missing parts become
0, and \texttt{has\_floor} is 1 when the string is present; % src: data_prep.py:clean (step 2)
(2)~land area: ``sot'' and spaces are removed, a comma becomes a decimal
point, and a missing value becomes 0; % src: data_prep.py:clean (step 3)
(3)~price, area, rooms, latitude, longitude and views are converted to
numbers after removing ``m\textsuperscript{2}'' and spaces and replacing
commas by points, unparsable values becoming missing. % src: data_prep.py:clean (step 4)
There is no currency conversion: every raw row is priced in AZN % src: results/report_facts.csv [raw_rows_currency_not_azn = 0]
and the currency columns are dropped; % src: data_prep.py:LEAKAGE_COLUMNS
(4)~rows missing price, area or rooms are removed (9{,}412 rows, all % src: results/report_facts.csv [rows_removed_missing]
without a room count), % src: results/report_facts.csv [raw_rows_rooms_missing = rows_removed_missing]
then exact duplicate rows (4); % src: results/report_facts.csv [rows_removed_duplicates]
missing values in seven text columns (location, city, building type,
category, renovation, title deed, mortgage) become ``Unknown''; % src: data_prep.py:clean (step 5)
(5)~outliers: only prices of 10{,}000--5{,}000{,}000~AZN and areas of
10--2{,}000~m\textsuperscript{2} are kept, % src: data_prep.py:clean (step 6)
removing 45 rows. % src: results/report_facts.csv [rows_removed_outliers]
This leaves 91{,}314 listings. % src: results/report_facts.csv [rows_clean]
No statistic of the data is computed during cleaning.

\textit{Encoding.} \texttt{make\_features} first turns the five flag
columns into presence indicators, 1 if the column has a value and 0 if
it is missing; one-hot encoding with \texttt{drop\_first} would otherwise
drop a single-valued column entirely. % src: data_prep.py:make_features (presence flags)
It then one-hot encodes every other text column with
\texttt{pandas.get\_dummies(drop\_first=True)} and keeps the numeric
columns as they are. % src: data_prep.py:make_features
This gives 171 features: % src: results/report_facts.csv [n_features]
9 numeric, % src: results/report_facts.csv [n_numeric_features]
5 presence flags % src: results/report_facts.csv [n_flag_features]
and 157 indicators, % src: results/report_facts.csv [n_onehot_features]
139 of them for location names. % src: results/report_facts.csv [n_onehot_features, location 139]
The flags are 1 for 85.5\% of listings (\texttt{repair}), % src: results/report_facts.csv [flag_share_repair 0.8551]
77.8\% (\texttt{bill\_of\_sale}), % src: results/report_facts.csv [flag_share_bill_of_sale 0.7783]
34.1\% (\texttt{mortgage}), % src: results/report_facts.csv [flag_share_mortgage 0.3414]
8.2\% (\texttt{vip}) % src: results/report_facts.csv [flag_share_vip 0.0824]
and 3.0\% (\texttt{featured}). % src: results/report_facts.csv [flag_share_featured 0.0303]
\texttt{repair}, \texttt{bill\_of\_sale} and \texttt{mortgage} are 1
on exactly the rows where the renovation, title-deed and mortgage columns
say ``yes'', % src: results/report_facts.csv [flag_mismatch_repair, flag_mismatch_bill_of_sale, flag_mismatch_mortgage = 0]
so they repeat information the one-hot indicators already carry;
\texttt{vip} and \texttt{featured} add new information.

\textit{Split and labels.} Every stage (EDA, SVM tuning, benchmark and
error analysis) takes its split from one function,
\texttt{get\_splits} in \texttt{splits.py}. % src: splits.py:get_splits; tests/test_splits.py
It calls \texttt{train\_val\_test\_split} with a stratification label:
the tier computed from the median of \emph{all} prices. This label only
keeps the premium share equal across the three splits; no model is
trained or scored on it (Section~\ref{sec:limits}). % src: splits.py:get_splits (strata)
The split is seeded and stratified at both steps: 15\% for test, then
15/85 of the rest for validation. % src: data_prep.py:train_val_test_split
It yields 63{,}919 training, % src: results/report_facts.csv [n_train]
13{,}697 validation % src: results/report_facts.csv [n_val]
and 13{,}698 test listings. % src: results/report_facts.csv [n_test]
\texttt{get\_splits} then recomputes the labels the models use with
\texttt{make\_tier\_label} on the training prices only: the threshold
is the median \emph{training} price, 215{,}000~AZN, % src: results/report_facts.csv [tier_threshold_train_median]
a listing is premium (1) if its price is at least the threshold, and % src: data_prep.py:make_tier_label
validation and test use the same threshold. % src: splits.py:get_splits
The classes are balanced: 50.3\% premium in each split. % src: results/report_facts.csv [premium_share_train 0.5027, premium_share_val 0.5027, premium_share_test 0.5027]
For the SVMs, \texttt{standardize} rescales every feature with the
training mean and standard deviation (a constant column is left
unscaled). % src: data_prep.py:standardize

\begin{table}[t]
  \centering
  \caption{Dataset summary. Features: numeric + presence flags + one-hot indicators.}
  \label{tab:data}
  \begin{tabular}{@{}lr@{}}
    \toprule
    Property & Value \\
    \midrule
    Raw rows / columns          & 100{,}775 / 51 \\ % src: results/report_facts.csv [raw_rows, raw_columns]
    Rows after cleaning         & 91{,}314 \\ % src: results/report_facts.csv [rows_clean]
    Features (see caption)      & 171 (9 + 5 + 157) \\ % src: results/report_facts.csv [n_features, n_numeric_features, n_flag_features, n_onehot_features]
    Train / val / test rows     & 63{,}919 / 13{,}697 / 13{,}698 \\ % src: results/report_facts.csv [n_train, n_val, n_test]
    Tier threshold (train median) & 215{,}000 AZN \\ % src: results/report_facts.csv [tier_threshold_train_median]
    Premium \%, train / val / test & 50.3 / 50.3 / 50.3 \\ % src: results/report_facts.csv [premium_share_train, _val, _test]
    \bottomrule
  \end{tabular}
\end{table}

% =====================================================================
\section{Exploratory Data Analysis}
\label{sec:eda}
We explored the training split only (63{,}919 listings), % src: results/eda_summary.csv [n_train]
so nothing seen here can leak validation or test prices into a choice. % src: eda.py:load_training_rows
Prices are strongly right-skewed (Fig.~\ref{fig:eda-price}): the mean is
271{,}961~AZN against a median of 215{,}000, % src: results/eda_summary.csv [price_mean 271961.3, price_median 215000]
and the skewness is 6.00. % src: results/eda_summary.csv [price_skew 6.0001]
The logarithm removes most of it (skewness 0.28), % src: results/eda_summary.csv [log_price_skew 0.2794]
which is why Task~A predicts $\log(\text{price})$: squared errors on the
log scale weigh a 10\% error equally at every price level. The tiers have
31{,}786 standard and 32{,}133 premium training listings. % src: results/eda_summary.csv [n_standard, n_premium]

\begin{figure*}[t]
  \centering
  \includegraphics[width=\textwidth]{figures/eda_price.pdf}
  \caption{Training split. Left: price, without the 640 listings above % src: results/eda_summary.csv [rows_above_p99]
  the 99th percentile (1{,}230{,}740~AZN). % src: results/eda_summary.csv [price_p99]
  Middle: $\log(\text{price})$, the Task~A target. Dashed: the tier
  threshold. Right: Task~B class balance.}
  \label{fig:eda-price}
\end{figure*}

Floor area is the strongest single signal (Fig.~\ref{fig:eda-features}):
log area and log price correlate at 0.73. % src: results/eda_summary.csv [corr_log_area_log_price 0.7313]
The median price rises with the number of rooms, from 105{,}500~AZN for % src: results/eda_summary.csv [median_price_rooms_1]
one room to 510{,}000~AZN for six or more, % src: results/eda_summary.csv [median_price_rooms_6_plus]
but the price ranges of neighbouring room counts overlap widely. On the
map, the most expensive listings form a dense cluster in central Baku, at
the south-west of the plotted area, and prices fall towards the outer
settlements. Location therefore matters beyond size, which the models
confirm (Section~\ref{sec:results}).

\begin{figure*}[t]
  \centering
  \includegraphics[width=\textwidth]{figures/eda_features.pdf}
  \caption{Training split. Left: floor area against price (log-log,
  darker = more listings). Middle: price by number of rooms (boxes:
  quartiles; whiskers: 1.5 IQR). Right: listing locations coloured by
  $\log(\text{price})$; the map shows the central 98\% of latitudes and % src: eda.py:MAP_QUANTILES (0.01, 0.99)
  longitudes, 61{,}593 of the 63{,}919 listings.} % src: results/eda_summary.csv [map_rows_shown, n_train]
  \label{fig:eda-features}
\end{figure*}

% =====================================================================
\section{Method}
\subsection{Decision tree (from scratch)}
\textit{Structure.} Each \texttt{Node} stores a feature index and
threshold (internal nodes), its two children, a leaf value, and the number
of training rows and the impurity at the node. % src: decision_tree.py:Node
\texttt{fit} records the classes present in $y$ and grows the tree
recursively from the root at depth 0. % src: decision_tree.py:DecisionTree.fit
The \texttt{random\_state} argument only seeds NumPy's global generator;
the algorithm itself uses no randomness. % src: decision_tree.py:DecisionTree.fit

\textit{Growth and stopping.} \texttt{\_grow} first stores the node's row
count and impurity. The node becomes a leaf if it is pure, has fewer than
\texttt{min\_samples\_split} rows, or has reached \texttt{max\_depth}; % src: decision_tree.py:DecisionTree._grow
otherwise it searches for the best split, and becomes a leaf if none is
valid. A valid split leaves at least \texttt{min\_samples\_leaf} rows in
each child % src: decision_tree.py:DecisionTree._best_split
and decreases the impurity by at least \texttt{min\_impurity\_decrease},
measured as~(\ref{eq:gain}) at the node (scikit-learn additionally
weights this test by the node's share of all rows). % src: decision_tree.py:DecisionTree._best_split
The default of this fourth parameter is 0, so a split with zero gain is
still accepted. % src: decision_tree.py:DecisionTree.__init__, _best_split (max_decrease >= min_impurity_decrease)
Rows with $x_j \le t$ go left, the others right.

\textit{Split search.} For every feature, the candidate thresholds are the
\emph{observed} distinct values of that feature in the node (not
midpoints between them, as in scikit-learn). For each candidate the
children's impurities are recomputed from scratch, and the split with the
largest decrease~(\ref{eq:gain}) wins; on a tie the first one found, i.e.\
the lowest feature index and smallest threshold, is kept. % src: decision_tree.py:DecisionTree._best_split
A node with $m$ rows thus costs $O(m)$ per candidate threshold, which
makes the search slow (Section~\ref{sec:results}).

\textit{Criteria.} For classification the impurity is Gini or entropy
with $\log_2$, as in~(\ref{eq:imp}). For regression it is the variance of
the node's targets (\texttt{np.var}), i.e.\ $I_{\text{MSE}}$, whatever
the \texttt{criterion} string. % src: decision_tree.py:DecisionTree._impurity

\textit{Prediction.} A regression leaf stores the mean target, a
classification leaf the class frequencies of its rows. % src: decision_tree.py:DecisionTree._calculate_leaf_value
A sample follows the thresholds to a leaf; \texttt{predict} returns the
mean, or the class with the highest frequency (a tie goes to the first
class, standard), and \texttt{predict\_proba} returns the frequencies. % src: decision_tree.py:DecisionTree.predict, predict_proba
We use the premium frequency as the ROC score. % src: benchmark.py:_fit_and_score ("proba")

\textit{Feature importance.} \texttt{tree\_feature\_importances} (in
\texttt{evaluate.py}) walks the fitted tree and adds, for every internal node, its weighted
impurity decrease $n_t I_t - n_L I_L - n_R I_R$ to the node's split
feature, then normalises the sum to~1, using the stored node
statistics. This is scikit-learn's mean decrease in impurity. % src: evaluate.py:tree_feature_importances

\textit{Correctness check.} \texttt{tests/test\_decision\_tree.py} fits
our tree and scikit-learn's with the same settings (\texttt{max\_depth=2})
on two six-row toy datasets and requires identical class predictions and
regression predictions equal to five decimals. % src: tests/test_decision_tree.py:test_classification_match, test_regression_match
Section~\ref{sec:results} repeats the comparison on the full data.

\subsection{Support-vector machine (from scratch)}
\textit{Objective and labels.} \texttt{PegasosSVM} minimises~(\ref{eq:svm})
with an unregularised bias $b$. % src: svm.py:PegasosSVM._record_objective, _fit_features
The two training classes are sorted and mapped to $-1$ and $+1$, so
standard is $-1$ and premium $+1$; \texttt{predict} returns premium when
the score $\mathbf{w}^\top\mathbf{z} + b \ge 0$. % src: svm.py:PegasosSVM._to_pm1, predict

\textit{Update as implemented.} Training starts from
$\mathbf{w} = \mathbf{0}$, $b = 0$ and makes one update per example; it is
single-example, not mini-batch, Pegasos. % src: svm.py:PegasosSVM._fit_features
With \texttt{n\_epochs} set, as in all our runs, the data is visited in a
new random order in each epoch, for $T = \texttt{n\_epochs}\cdot n$
steps. % src: svm.py:PegasosSVM._fit_features (n_steps, rng.permutation)
At step $t$ the code computes the margin $y_i(\mathbf{w}^\top\mathbf{z}_i + b)$
with the current parameters, sets $\eta_t = 1/(\lambda t^{p})$ with
$p = 1$ (the default, used throughout), shrinks
$\mathbf{w} \leftarrow (1 - \eta_t\lambda)\mathbf{w}$ and, if the margin
is below~1, adds $\eta_t y_i \mathbf{z}_i$ to $\mathbf{w}$ and
$\eta_t y_i$ to $b$. % src: svm.py:PegasosSVM._fit_features
Finally $\mathbf{w}$ (not $b$) is projected onto the ball of radius
$1/\sqrt{\lambda}$. % src: svm.py:PegasosSVM._fit_features (radius = 1/sqrt(lambda_))
This is update~(\ref{eq:pegasos}) plus a bias that takes the same
sub-gradient step but is never shrunk or projected. The
\texttt{batch\_size} parameter only splits the Fourier mapping and
prediction into blocks to bound memory. % src: svm.py:RandomFourierFeatures.transform, PegasosSVM.decision_function

\textit{Diagnostics.} The model records the objective~(\ref{eq:svm}) and
the mean hinge loss on all training rows at step 0, every
\texttt{history\_every} steps and at the last step; % src: svm.py:PegasosSVM._record_objective, _fit_features
we set \texttt{history\_every} to $n$, one checkpoint per epoch. % src: svm_tuning.py:run_tuning, benchmark.py:_StandardizedPegasos.fit
After training it stores $\lVert\mathbf{w}\rVert$, the margin width
$2/\lVert\mathbf{w}\rVert$, and as support vectors the training rows with
$y_i(\mathbf{w}^\top\mathbf{z}_i + b) \le 1 + 10^{-10}$, i.e.\ on or
inside the margin, including misclassified rows. This margin-based
definition differs from the non-zero dual coefficients of a QP
solver. % src: svm.py:PegasosSVM._fit_features (support_)

\textit{Fourier features.} The RFF variant uses the class
\texttt{RandomFourierFeatures}, which draws $W$ and $\mathbf{a}$ as in
Section~\ref{sec:background} from a seeded generator, once, when the model
is fit, and maps every input with $\mathbf{z}(\mathbf{x})=\sqrt{2/D}\cos(\mathbf{x}W+\mathbf{a})$; % src: svm.py:RandomFourierFeatures.fit, transform
Pegasos then learns $\mathbf{w} \in \mathbb{R}^D$ and $b$. We use
$D = 500$. % src: results/report_facts.csv [svm_rff_n_components]
RFF was chosen as the non-linear option for memory: an exact float64
kernel matrix on our training split would need 32.7~GB, % src: results/report_facts.csv [kernel_matrix_gb 32.69]
the 500-dimensional Fourier matrix 256~MB. % src: results/report_facts.csv [rff_matrix_mb 255.7]
The RFF model did not, however, win the validation comparison: the final
SVM is linear (Section~\ref{sec:tuning}). Both variants are sweepable
with the same code, \texttt{svm\_experiment.sweep\_svm} for RFF. % src: svm_experiment.py:sweep_svm

% =====================================================================
\section{Experimental Setup}
\label{sec:setup}

\subsection{Targets, metrics and scaling}
Task~A predicts $\log(\text{price})$; we report RMSE, MAE and $R^2$ on
the log scale and, after $\exp(\cdot)$, in AZN. For Task~B we report
precision, recall and F1 for the premium class and ROC-AUC computed from
tree leaf probabilities or SVM decision values. All metrics are our own
NumPy code, tested against \texttt{sklearn.metrics}. % src: evaluate.py, tests/test_evaluate.py
The SVMs receive standardized features; the trees use the unscaled
features. % src: benchmark.py:_StandardizedPegasos, _task_a_models, _task_b_models

\subsection{Cross-validation and the held-out test}
No choice uses the test split. Cross-validation uses a stratified sample
of 20{,}000 % src: results/report_facts.csv [cv_rows]
of the 77{,}616 training and validation rows % src: results/report_facts.csv [cv_pool_rows]
and 5-fold stratified K-fold. % src: results/report_facts.csv [cv_folds]
Inside every fold, the tier threshold and the SVM scaler are recomputed
from that fold's training rows only, and every model sees the same folds
and features. We report the mean $\pm$ sample standard deviation over
folds. The subsample keeps the run practical: our tree takes about 10~s
per fit at 16{,}000 rows % src: results/benchmark_cv_task_b.csv [Tree gini (ours), Fit s mean 10.27]
and 87~s on the full training split. % src: results/benchmark_test_task_b.csv [Tree gini (ours), Fit s 86.75]
For the held-out test, each model is fit once on the full training split
and scored once on the test split.

\subsection{Matched scikit-learn baselines}
Our trees are compared with \texttt{DecisionTreeRegressor} and
\texttt{DecisionTreeClassifier} (Gini and entropy) with the same
\texttt{max\_depth}, \texttt{min\_samples\_split},
\texttt{min\_samples\_leaf} and seed. Linear Pegasos is compared with
\texttt{SGDClassifier} (hinge loss, L2 penalty, $\alpha = \lambda$, the
same number of epochs), which minimises the same objective~(\ref{eq:svm}),
and with \texttt{SVC(kernel="linear")}. RFF Pegasos is compared with the
same random-feature construction (\texttt{RBFSampler}, same $\gamma$ and
dimension) followed by \texttt{SGDClassifier}, and with
\texttt{SVC(kernel="rbf")} at the same $\gamma$. % src: benchmark.py:_svm_models, _sgd_pipeline, _SubsampledSVC
For \texttt{SVC} we set $C = 1/(n\lambda)$, which makes its objective a
rescaling of~(\ref{eq:svm}); $n$ is the number of rows SVC actually sees. % src: benchmark.py:_SubsampledSVC.fit
Exact-kernel SVC grows roughly quadratically with $n$, so it is fit on at
most 10{,}000 stratified training rows; % src: results/report_facts.csv [svc_max_rows]
all other models use every training row.

\subsection{Timing and seeds}
We record wall-clock fit and prediction time for every model on one
machine (Apple M4), and the fit time of both trees and both linear SVM
solvers at 2{,}000, 4{,}000, 8{,}000 and 16{,}000 training rows. % src: results/benchmark_scaling.csv [n_train]
Every random choice (split, subsamples, folds, models, random features)
uses the seed 42. % src: results/report_facts.csv [seed]

% =====================================================================
\section{Hyperparameter Tuning}
\label{sec:tuning}
\subsection{SVM: regularisation and kernel width}
Both SVMs were tuned on train $\to$ validation, fitting on all 63{,}919 % src: results/report_facts.csv [n_train]
training rows and scoring on the 13{,}697 validation rows. % src: results/report_facts.csv [n_val]
The grid has 7 values of $\lambda$, % src: results/report_facts.csv [svm_grid_n_lambdas]
from $10^{-7}$ to $10^{-1}$, % src: results/report_facts.csv [svm_grid_lambda_min, svm_grid_lambda_max]
for the linear kernel and for RFF with $D = 500$ % src: results/report_facts.csv [svm_rff_n_components]
and $\gamma \in \{0.1, 1, 10\} \times 1/171$, % src: results/report_facts.csv [svm_grid_gamma_multipliers, n_features]
where $1/171$ (one over the number of features) is the usual
\emph{scale} default for standardized inputs; each model trains for
20 epochs. % src: results/report_facts.csv [svm_linear_n_epochs]
The selection metric is validation F1; ROC-AUC is recorded too. % src: svm_tuning.py:_f1_recording_auc

Fig.~\ref{fig:tuning} shows the sweep. The linear SVM peaks at
$\lambda = 10^{-3}$ with F1 0.904. % src: results/svm_tuning/sweep_table.csv [linear, 0.001, val_F1 0.9037]
Smaller $\lambda$ lets $\lVert\mathbf{w}\rVert$ grow to 2{,}716 at
$10^{-7}$ % src: results/svm_tuning/sweep_table.csv [linear, 1e-07, weight_norm 2715.7]
and F1 fall to 0.866; % src: results/svm_tuning/sweep_table.csv [linear, 1e-07, val_F1 0.8664]
larger $\lambda$ widens the margin (1.71 at $10^{-1}$, against 0.45 at % src: results/svm_tuning/sweep_table.csv [linear, 0.1, margin_width 1.7097; 0.001, 0.4451]
$10^{-3}$) and puts 32{,}036 training rows on or inside it, % src: results/svm_tuning/sweep_table.csv [linear, 0.1, support_count 32036]
underfitting at F1 0.860. % src: results/svm_tuning/sweep_table.csv [linear, 0.1, val_F1 0.8600]
RFF does best with small $\lambda$ and the widest kernel (smallest
$\gamma$): its best F1 is 0.903, at $\gamma = 0.1 \times 1/171$ and % src: results/svm_tuning/sweep_table.csv [rff, 1e-06, gamma_x_scale 0.1, val_F1 0.9027]
$\lambda = 10^{-6}$, while $\gamma = 10 \times 1/171$ never exceeds 0.817. % src: results/svm_tuning/sweep_table.csv [rff, gamma_x_scale 10, max val_F1 0.8173]
At $\lambda = 10^{-1}$ every RFF model scores F1 0.669, % src: results/svm_tuning/sweep_table.csv [rff, 0.1, val_F1 0.6691 for all three gammas]
the F1 of labelling every validation listing premium. % src: 2 * 0.5027 / (1 + 0.5027) = 0.6691 from results/report_facts.csv [premium_share_val]

\begin{figure}[t]
  \centering
  \includegraphics[width=\linewidth]{figures/svm_validation_f1.pdf}
  \caption{SVM tuning: validation F1 against $\lambda$ for the linear
  kernel and RFF at three kernel widths ($\textit{scale} = 1/171$). Stars:
  the configuration chosen for each kernel; filled: the final model.}
  \label{fig:tuning}
\end{figure}

\textit{Selection rule and stability check.} % every number in this paragraph: results/svm_tuning/stability_check.csv
We do not simply take the highest F1. Within each kernel, among
configurations whose validation F1 is within 0.005 of the best, we take % src: results/svm_tuning/stability_check.csv [f1_tolerance]
the \emph{largest}~$\lambda$. % src: svm_tuning.py:_select_stable
Pegasos uses the step size $1/(\lambda t)$ and its guarantee after $T$
updates scales as $1/(\lambda T)$, so it needs $T \gg 1/\lambda$, and a
small $\lambda$ that wins on the full training split can be
under-trained on less data. To test this, each kernel's
highest-validation-F1 configuration and its chosen configuration are
re-evaluated with the benchmark's cross-validation design: % src: svm_tuning.py:stability_check
20{,}000 training and validation rows, % src: results/svm_tuning/stability_check.csv [cv_rows]
5 folds, % src: results/svm_tuning/stability_check.csv [cv_folds]
16{,}000 training rows per fold % src: results/svm_tuning/stability_check.csv [fit_rows_per_fold]
and 20 epochs, i.e.\ 320{,}000 updates per fit. % src: results/svm_tuning/stability_check.csv [n_epochs, updates_per_fit]
For the linear kernel both configurations are $\lambda = 10^{-3}$, % src: results/svm_tuning/stability_check.csv [linear, lambda_ 0.001 in both rows]
and its validation F1 of 0.904 holds up in cross-validation at % src: results/svm_tuning/stability_check.csv [linear, val_F1 0.9037]
$0.897 \pm 0.003$. % src: results/svm_tuning/stability_check.csv [linear, cv_F1_mean 0.8970, cv_F1_std 0.0031]
For RFF both are $\lambda = 10^{-6}$ with $\gamma = 0.1 \times \textit{scale}$: % src: results/svm_tuning/stability_check.csv [rff, lambda_ 1e-06, gamma_x_scale 0.1 in both rows]
no configuration with a larger $\lambda$ came within 0.005 of its % src: results/svm_tuning/stability_check.csv [f1_tolerance]
validation F1 of 0.903, so the rule did not change the choice. % src: results/svm_tuning/stability_check.csv [rff, val_F1 0.9027]
This configuration needs $1/\lambda = 10^6$ updates % src: results/svm_tuning/stability_check.csv [rff, one_over_lambda 1e6]
but gets 320{,}000 per fold, % src: results/svm_tuning/stability_check.csv [rff, updates_per_fit]
and its cross-validation F1 falls to $0.763 \pm 0.095$. % src: results/svm_tuning/stability_check.csv [rff, cv_F1_mean 0.7632, cv_F1_std 0.0954]
The rule only guards against this when a larger $\lambda$ is nearly as
good on validation, which was not the case for RFF on this split.

\textit{Final SVM.} The final SVM is the kernel whose chosen configuration
has the higher validation F1: linear, % src: results/report_facts.csv [svm_final_kernel]
0.904 against 0.903 for RFF, % src: results/report_facts.csv [svm_linear_val_f1 0.9037, svm_rff_val_f1 0.9027]
although RFF has the marginally higher validation ROC-AUC (0.9647 % src: results/report_facts.csv [svm_rff_val_auc]
against 0.9637). % src: results/report_facts.csv [svm_linear_val_auc]
Both SVMs are benchmarked.

\subsection{Decision tree}
Both tasks use one fixed tree configuration, set before any experiment
and not tuned: maximum depth 10, % src: results/report_facts.csv [tree_task_a_max_depth, tree_task_b_max_depth]
at least 10 rows to split a node, % src: results/report_facts.csv [tree_task_a_min_samples_split, tree_task_b_min_samples_split]
at least 5 rows per leaf, % src: results/report_facts.csv [tree_task_a_min_samples_leaf, tree_task_b_min_samples_leaf]
and no minimum impurity decrease. % src: benchmark.py:FALLBACK_CONFIGS, _tree_params (not set) -> decision_tree.py:DecisionTree.__init__ default 0.0
These are the fixed values in \texttt{benchmark.py}; the pipeline runs no
tree sweep, so nothing overrides them. % src: results/report_facts.csv [tree_task_a_source, tree_task_b_source]; run_all.py:_tree_tuning
For Task~B we compare the two criteria in cross-validation: Gini and
entropy reach the same F1, % src: results/benchmark_cv_task_b.csv [Tree gini (ours), F1 0.9069 ± 0.0060; Tree entropy (ours), F1 0.9072 ± 0.0035]
$0.907 \pm 0.006$ and $0.907 \pm 0.004$, % src: results/benchmark_cv_task_b.csv [Tree gini (ours) / Tree entropy (ours), F1]
while entropy ranks better (ROC-AUC $0.962 \pm 0.004$ % src: results/benchmark_cv_task_b.csv [Tree entropy (ours), ROC-AUC 0.9616 ± 0.0035]
against $0.951 \pm 0.005$) % src: results/benchmark_cv_task_b.csv [Tree gini (ours), ROC-AUC 0.9512 ± 0.0046]
and fits faster (9.1~s against 10.3~s per fold). % src: results/benchmark_cv_task_b.csv [Fit s, Tree entropy (ours) 9.0930, Tree gini (ours) 10.2740]
We report the Gini tree as our main Task~B tree because that choice was
fixed in advance; on F1 the data gives no reason to prefer either.

% =====================================================================
\section{Results}
\label{sec:results}

\begin{table*}[t]
  \centering
  \small
  \caption{Headline results on the held-out test set, with the 5-fold CV
  mean $\pm$ std where marked ``CV''. Task~A errors are on the log scale
  unless marked AZN; Task~B metrics are for the premium class. n/a: the
  model is for the other task. $^{*}$Fit on at most 10{,}000 training
  rows. All metrics: Tables~\ref{tab:full-a} and~\ref{tab:full-b}.}
  \label{tab:main}
  \begin{tabular}{lrrrrrrr}
\toprule
 & \multicolumn{4}{c}{Task A: log price} & \multicolumn{3}{c}{Task B: premium tier} \\
\cmidrule(lr){2-5}\cmidrule(lr){6-8}
Model & RMSE (log) & $R^2$ (log) & $R^2$ (log), CV & $R^2$ (AZN) & F1 & F1, CV & ROC-AUC \\
\midrule
Tree (ours) & 0.227 & 0.870 & 0.838 $\pm$ 0.008 & 0.764 & n/a & n/a & n/a \\
Tree (sklearn) & 0.227 & 0.870 & 0.839 $\pm$ 0.008 & 0.753 & n/a & n/a & n/a \\
Tree gini (ours) & n/a & n/a & n/a & n/a & 0.926 & 0.907 $\pm$ 0.006 & 0.971 \\
Tree gini (sklearn) & n/a & n/a & n/a & n/a & 0.926 & 0.907 $\pm$ 0.006 & 0.971 \\
Tree entropy (ours) & n/a & n/a & n/a & n/a & 0.920 & 0.907 $\pm$ 0.004 & 0.975 \\
Tree entropy (sklearn) & n/a & n/a & n/a & n/a & 0.920 & 0.909 $\pm$ 0.004 & 0.975 \\
Pegasos linear (ours) & n/a & n/a & n/a & n/a & 0.904 & 0.897 $\pm$ 0.003 & 0.962 \\
SGD hinge linear (sklearn) & n/a & n/a & n/a & n/a & 0.905 & 0.894 $\pm$ 0.004 & 0.962 \\
SVC linear (sklearn)$^{*}$ & n/a & n/a & n/a & n/a & 0.904 & 0.903 $\pm$ 0.002 & 0.960 \\
Pegasos RFF (ours) & n/a & n/a & n/a & n/a & 0.900 & 0.763 $\pm$ 0.095 & 0.963 \\
SGD hinge RFF (sklearn) & n/a & n/a & n/a & n/a & 0.887 & 0.829 $\pm$ 0.047 & 0.964 \\
SVC RBF (sklearn)$^{*}$ & n/a & n/a & n/a & n/a & 0.907 & 0.903 $\pm$ 0.001 & 0.965 \\
\bottomrule
\end{tabular}

\end{table*}

Table~\ref{tab:main} lists every model. On Task~A our tree reaches a test
$R^2$ of 0.853 on log price % src: results/benchmark_test_task_a.csv [Tree (ours), R2 log 0.8531]
(0.753 in AZN), % src: results/benchmark_test_task_a.csv [Tree (ours), R2 AZN 0.7531]
against 0.852 for the scikit-learn tree; % src: results/benchmark_test_task_a.csv [Tree (sklearn), R2 log 0.8522]
in cross-validation the two score $0.831 \pm 0.006$ % src: results/benchmark_cv_task_a.csv [Tree (ours), R2 log 0.8307 ± 0.0056]
and $0.831 \pm 0.005$. % src: results/benchmark_cv_task_a.csv [Tree (sklearn), R2 log 0.8310 ± 0.0049]
The gap between log scale and AZN comes from the most expensive
listings (Section~\ref{sec:errors}).

On Task~B the trees lead. Our Gini tree reaches test F1 0.926 % src: results/benchmark_test_task_b.csv [Tree gini (ours), F1 0.9258]
and ROC-AUC 0.971, % src: results/benchmark_test_task_b.csv [Tree gini (ours), ROC-AUC 0.9705]
equal to scikit-learn's tree (F1 0.926). % src: results/benchmark_test_task_b.csv [Tree gini (sklearn), F1 0.9256]
The final linear Pegasos SVM reaches F1 0.904 % src: results/benchmark_test_task_b.csv [Pegasos linear (ours), F1 0.9040]
and ROC-AUC 0.962, % src: results/benchmark_test_task_b.csv [Pegasos linear (ours), ROC-AUC 0.9620]
level with its baselines: SGD 0.905 % src: results/benchmark_test_task_b.csv [SGD hinge linear (sklearn), F1 0.9052]
and linear SVC 0.904. % src: results/benchmark_test_task_b.csv [SVC linear (sklearn, max 10k rows), F1 0.9045]
RFF Pegasos gets F1 0.900, % src: results/benchmark_test_task_b.csv [Pegasos RFF (ours), F1 0.9003]
above RFF SGD (0.887) % src: results/benchmark_test_task_b.csv [SGD hinge RFF (sklearn), F1 0.8872]
and below exact RBF SVC (0.907). % src: results/benchmark_test_task_b.csv [SVC RBF (sklearn, max 10k rows), F1 0.9068]
Each of our models is within 0.007 F1 of its best matched baseline; % src: results/error_analysis.md [Findings, bullet 1]
the largest gap is RFF Pegasos against the exact RBF SVC. % src: results/error_analysis.md [Findings, bullet 1: largest gap Pegasos RFF (ours), -0.0065]
Trained on the full split, RFF Pegasos is competitive, but it is the least
stable model in cross-validation, as the stability check predicted
(Section~\ref{sec:tuning}): F1 $0.763 \pm 0.095$, % src: results/benchmark_cv_task_b.csv [Pegasos RFF (ours), F1 0.7632 ± 0.0954]
against $0.897 \pm 0.003$ for linear Pegasos. % src: results/benchmark_cv_task_b.csv [Pegasos linear (ours), F1 0.8970 ± 0.0031]
SGD on the same random features is also unstable ($0.829 \pm 0.047$), % src: results/benchmark_cv_task_b.csv [SGD hinge RFF (sklearn), F1 0.8286 ± 0.0472]
while the exact RBF SVC is not ($0.903 \pm 0.001$). % src: results/benchmark_cv_task_b.csv [SVC RBF (sklearn, max 10k rows), F1 0.9032 ± 0.0014]
Fig.~\ref{fig:roc} shows that the ROC curves of our models and their
baselines nearly coincide.

\begin{figure*}[t]
  \centering
  \includegraphics[width=\textwidth]{figures/analysis_roc_curves.pdf}
  \caption{Task~B ROC curves on the test set by model family, zoomed on
  false-positive rates up to 0.3. Solid: our models; dashed:
  scikit-learn. Legend numbers: ROC-AUC. SVC is fit on at most 10{,}000
  rows.}
  \label{fig:roc}
\end{figure*}

\textit{Our trees versus scikit-learn's.} The trees are practically the
same model. Both Task~B Gini trees have depth 10, % src: results/tree_comparison.csv [Task B gini, depth ours/sklearn]
464 and 462 leaves, % src: results/tree_comparison.csv [Task B gini, leaves ours/sklearn]
split first on floor area, and agree on 99.9\% of test labels. % src: results/tree_comparison.csv [Task B gini, test agreement % 99.89]
The root thresholds differ only by convention: 90.1~m\textsuperscript{2} % src: results/tree_comparison.csv [Task B gini, root ours]
for ours, an observed training value, and 90.15~m\textsuperscript{2} % src: results/tree_comparison.csv [Task B gini, root sklearn]
for scikit-learn, the midpoint between neighbouring values. The Task~A
trees both have 656 leaves; % src: results/tree_comparison.csv [Task A tree, leaves ours/sklearn]
their predictions correlate at 0.9989, % src: results/tree_comparison.csv [Task A tree, pred corr 0.99886]
differ by 0.0015 log units on average, % src: results/tree_comparison.csv [Task A tree, mean |diff| log 0.00148]
and 99.6\% agree within 0.01 log units (about 1\% of the price). % src: results/tree_comparison.csv [Task A tree, within 0.01 log % 99.63]
The feature importances correlate above 0.9999 for all three pairs % src: results/tree_comparison.csv [importance corr, min 0.99996]
(Fig.~\ref{fig:importance}). Floor area dominates both tasks, followed
by longitude and latitude, in line with the EDA.

\begin{figure*}[t]
  \centering
  \includegraphics[width=\textwidth]{figures/analysis_feature_importance.pdf}
  \caption{Feature importance (mean decrease in impurity, normalised) of
  our trees and scikit-learn's: the 15 most important features by our
  importance. Location names end in q.\ (settlement), m.\ (metro-station
  area) or r.\ (district).}
  \label{fig:importance}
\end{figure*}

\textit{Speed.} Matching the models' quality costs time for the tree. On
the full training split our Gini tree fits in 86.8~s % src: results/benchmark_test_task_b.csv [Tree gini (ours), Fit s 86.7520]
against 0.31~s for scikit-learn, % src: results/benchmark_test_task_b.csv [Tree gini (sklearn), Fit s 0.3128]
about 277 times slower. % src: 86.7520 / 0.3128 from results/benchmark_test_task_b.csv
Fig.~\ref{fig:scaling} shows the growth with $n$: from 2{,}000 to
16{,}000 rows our tree's fit time grows 15.0-fold % src: results/benchmark_scaling.csv [Tree gini (ours), 10.3910 / 0.6940]
for 8 times the data, scikit-learn's 10.4-fold, % src: results/benchmark_scaling.csv [Tree gini (sklearn), 0.06312 / 0.006055]
and the two linear SVM solvers 8.6-fold and 9.3-fold, close to % src: results/benchmark_scaling.csv [Pegasos linear 0.6344/0.07362 = 8.62; SGD 0.05857/0.006320 = 9.27]
linearly. Linear Pegasos fits the full training split in 2.7~s, % src: results/benchmark_test_task_b.csv [Pegasos linear (ours), Fit s 2.6653]
against 0.29~s for SGD. % src: results/benchmark_test_task_b.csv [SGD hinge linear (sklearn), Fit s 0.2875]

\begin{figure}[t]
  \centering
  \includegraphics[width=\linewidth]{figures/benchmark_scaling.pdf}
  \caption{Fit time against training size (log-log), Task~B, one fit per
  point. Slope: fitted exponent $k$ in time $\propto n^k$. Our tree grows
  faster than linearly; both SVM solvers grow close to linearly.} % src: results/benchmark_meta.json [scaling_slopes: 1.29, 1.12, 1.03, 1.06]
  \label{fig:scaling}
\end{figure}

% =====================================================================
\section{Error Analysis}
\label{sec:errors}
We analyse the test errors of our Task~A tree and, on Task~B, of our Gini
tree and the final linear SVM.

\textit{The tier boundary.} Task~B errors are roughly balanced between
the classes (precision 0.921, recall 0.930 for the tree; % src: results/benchmark_test_task_b.csv [Tree gini (ours), Precision 0.9212, Recall 0.9304]
Fig.~\ref{fig:confusion}) and concentrate near the threshold. For
listings within 5\% of 215{,}000~AZN the tree misclassifies 32.3\% % src: results/error_analysis.md [Task B: distance from the tier threshold, 0-5%, B error % tree 32.31]
and the SVM 42.8\%, % src: results/error_analysis.md [Task B: distance from the tier threshold, 0-5%, B error % SVM 42.81]
while beyond 40\% from the threshold they err on 1.3\% and 1.5\%. % src: results/error_analysis.md [Task B: distance from the tier threshold, >= 40%, 1.31 / 1.48]
The 5\% band alone holds 28\% of all tree errors. % src: results/error_analysis.md [Findings, bullet 5]
Most of the remaining Task~B error is therefore noise in listing prices
close to an arbitrary cut-off.

\begin{figure}[t]
  \centering
  \includegraphics[width=\linewidth]{figures/analysis_confusion_matrices.pdf}
  \caption{Task~B confusion matrices on the test set, our Gini tree and
  our final linear Pegasos SVM (counts and share of the true tier).}
  \label{fig:confusion}
\end{figure}

\textit{Price extremes.} On Task~A the median absolute percentage error
is 14.3\% in the cheapest decile % src: results/error_analysis.md [By true-price decile, D1, A MdAPE % 14.27]
and 19.6\% in the most expensive, % src: results/error_analysis.md [By true-price decile, D10, A MdAPE % 19.64]
against 9.3--12.5\% for deciles 2--9 % src: results/error_analysis.md [By true-price decile, D2-D9, A MdAPE % 9.29-12.49]
(Fig.~\ref{fig:deciles}). The errors are biased towards the middle: the
cheapest decile is over-predicted by a median 7.8\% % src: results/error_analysis.md [By true-price decile, D1, A median bias % 7.84]
and the most expensive under-predicted by 13.8\%. % src: results/error_analysis.md [By true-price decile, D10, A median bias % -13.81]
A tree of depth 10 predicts leaf means, so prices at the extremes are
pulled towards the average of their leaf. The effect is strongest for
the 138 most expensive listings (top 1\%, from 1{,}300{,}000~AZN): % src: results/error_analysis.md [Findings, bullet 4: n = 138, price >= 1,300,000 AZN]
median absolute error 33.9\% % src: results/error_analysis.md [Top 1% most expensive listings, A MdAPE % 33.88]
and a median under-prediction of 33.6\%. % src: results/error_analysis.md [Top 1% most expensive listings, A median bias % -33.60]
For Task~B the same listings are easy (tree error 0.0\%), % src: results/error_analysis.md [Top 1% most expensive listings, B error % tree 0.00]
being far from the threshold.

\begin{figure*}[t]
  \centering
  \includegraphics[width=\textwidth]{figures/analysis_error_by_price_decile.pdf}
  \caption{Errors by true-price decile (decile edges from training
  prices), test set. Left: Task~A median absolute percentage error of our
  tree. Right: Task~B error rate; the threshold lies between deciles 5
  and 6.}
  \label{fig:deciles}
\end{figure*}

\textit{Sparse districts.} Task~A error grows as a district's training
data shrinks. Grouping test listings by how many \emph{training} listings
their district has, the median absolute error is 11.4\% % src: results/error_analysis.md [By district size, >= 1000, A MdAPE % 11.35]
for districts with 1{,}000 or more training rows, % src: results/error_analysis.md [By district size, bin >= 1000]
13.1\% for 200--999, % src: results/error_analysis.md [By district size, 200-999, A MdAPE % 13.05]
18.6\% for 50--199 % src: results/error_analysis.md [By district size, 50-199, A MdAPE % 18.60]
and 20.0\% for 1--49. % src: results/error_analysis.md [By district size, 1-49, A MdAPE % 20.03]
Task~B shows no such trend: the tree's error rate is 6.9\% in the % src: results/error_analysis.md [By district size, >= 1000, B error % tree 6.91]
largest districts, 11.8\% for 50--199 % src: results/error_analysis.md [By district size, 50-199, B error % tree 11.80]
and 5.8\% for 1--49, whose listings are mostly cheap (15.7\% premium) % src: results/error_analysis.md [By district size, 1-49, B error % tree 5.76, B premium % 15.71]
and therefore far from the threshold.
One test listing comes from a district with no training rows. % src: results/error_analysis.md [By district size, 0 (unseen in train), n test 1]

\textit{Categories.} Houses and villas are the hardest large group, with
a median absolute error of 19.1\% % src: results/error_analysis.md [By category, House / villa, A MdAPE % 19.09]
against 11.3\% for new-build % src: results/error_analysis.md [By category, New build, A MdAPE % 11.33]
and 10.3\% for old-build apartments. % src: results/error_analysis.md [By category, Old build, A MdAPE % 10.25]
Offices score 34.4\%, % src: results/error_analysis.md [By category, Office, A MdAPE % 34.36]
but there are only 25 of them in the test set. % src: results/error_analysis.md [By category, Office, n test 25]

\textit{Worst listings.} Nine of the ten worst Task~A listings are houses % src: results/error_analysis.md [10 worst Task A listings, category: 9 x House / villa]
or villas whose price differs from our prediction by a factor of 4.8 % src: results/error_analysis/worst_task_a_listings.csv [houses: max(pred/true, true/pred) 4.77 to 8.69]
to 8.7 in either direction. Some look like data-entry errors rather than model failures,
such as a seven-room house with 11~m\textsuperscript{2} of floor area on % src: results/error_analysis.md [10 worst Task A listings, clean_index 90432, rooms 7, area_m2 11]
310 sot, listed at 280{,}000~AZN. % src: results/error_analysis.md [10 worst Task A listings, clean_index 90432, land_sot 310, price 280,000]
No model can price such listings, and they inflate RMSE in AZN.

\textit{What the SVM uses.} The final linear SVM has
$\lVert\mathbf{w}\rVert = 4.49$, % src: results/error_analysis.md [Final linear SVM, ||w|| 4.4939]
a margin $2/\lVert\mathbf{w}\rVert = 0.445$ in standardized units, % src: results/error_analysis.md [Final linear SVM, margin width 0.4451]
and 17{,}946 of the 63{,}919 training points on or inside the margin. % src: results/error_analysis.md [Final linear SVM, support vectors 17,946 of 63,919]
Floor area carries by far the largest weight towards premium; being in
an outlying settlement of Baku carries the largest weights towards
standard (all weights: \url{results/error_analysis/svm_linear_weights.csv}). % src: results/error_analysis/svm_linear_weights.csv [Sahə 3.89, next largest has_floor 0.74; most negative Masazır q., Binə q., Binəqədi q.]

% =====================================================================
\section{Limitations}
\label{sec:limits}
\begin{itemize}
  \item \textit{The trees were not tuned.} Both tasks use the fixed
  configuration in \texttt{src/benchmark.py} % src: results/report_facts.csv [tree_task_a_source, tree_task_b_source]
  (maximum depth 10, % src: results/report_facts.csv [tree_task_a_max_depth, tree_task_b_max_depth]
  minimum 10 samples to split, % src: results/report_facts.csv [tree_task_a_min_samples_split, tree_task_b_min_samples_split]
  minimum 5 per leaf), % src: results/report_facts.csv [tree_task_a_min_samples_leaf, tree_task_b_min_samples_leaf]
  because the pipeline contains no tree sweep on the real data, while the
  SVMs were tuned on validation. The trees' results are therefore not their best, and the
  tuning effort between the two model families is unequal.
  \item \textit{Cross-validation uses a subsample.} CV uses 20{,}000 % src: results/report_facts.csv [cv_rows]
  of the 77{,}616 available rows, % src: results/report_facts.csv [cv_pool_rows]
  so each fold trains on 16{,}000 rows, % src: results/report_facts.csv [cv_fit_rows_per_fold]
  a quarter of the full training split. CV scores are
  correspondingly lower than the test scores (Task~A $R^2$ $0.831$ % src: results/benchmark_cv_task_a.csv [Tree (ours), R2 log mean 0.8307]
  against $0.853$) % src: results/benchmark_test_task_a.csv [Tree (ours), R2 log 0.8531]
  and describe the variability of a smaller-data model; for RFF Pegasos
  the smaller folds are enough to break training (Section~\ref{sec:tuning}).
  \item \textit{One-hot categories come from all rows.} The encoding is
  built before the split, so categories that occur only in validation or
  test still get a column. Two such columns are all-zero in training. % src: results/report_facts.csv [onehot_all_zero_in_train = 2]
  No statistic leaks, and the models never use these columns, but a
  strict pipeline would fit the encoder on training rows only.
  \item \textit{Stratification on a full-data label.} \texttt{get\_splits}
  balances the three splits with a tier computed from the median of
  \emph{all} prices, test prices included, so which listings land in the
  test split depends on that median. % src: splits.py:get_splits
  No model is trained or scored on this label: the labels the models use
  come from the training median. On this data both medians are
  215{,}000~AZN, % src: results/report_facts.csv [full_data_median, tier_threshold_train_median]
  so the two labelings differ on 0 listings. % src: results/report_facts.csv [labels_differ_full_vs_train_median]
  \item \textit{Our tree is slow.} It is about 277 times slower than % src: 86.7520 / 0.3128 from results/benchmark_test_task_b.csv
  scikit-learn on the full training split and grows faster than
  linearly with $n$ (Fig.~\ref{fig:scaling}). For every candidate
  threshold it recomputes the impurity from scratch instead of sweeping
  sorted values with running sums. This forced the CV subsample and left
  no budget to tune the tree.
  \item \textit{Two fields are filled with 0.} For the 15{,}338 listings % src: results/report_facts.csv [rows_no_floor_filled_0]
  without a floor (15{,}188 houses and 150 offices), % src: results/report_facts.csv [rows_no_floor_house, rows_no_floor_office]
  floor and building height are set to 0, with a \texttt{has\_floor}
  flag; for the 76{,}126 listings without a land area, % src: results/report_facts.csv [rows_land_area_filled_0]
  the land area is set to 0. Here 0 means ``not applicable'', not a
  measured zero. The trees can separate the two through the flag and the
  category, but the linear SVM treats 0 as a real value.
\end{itemize}

% =====================================================================
\section{Conclusion}
Our from-scratch decision tree and Pegasos SVM reach the accuracy of
matched scikit-learn models on bina.az listings: each is within 0.007 F1 % src: results/error_analysis.md [Findings, bullet 1]
of its best baseline, and our trees reproduce scikit-learn's almost split
for split. The tree is the better model on both tasks, with test $R^2$ 0.853 % src: results/benchmark_test_task_a.csv [Tree (ours), R2 log 0.8531]
on log price and F1 0.926 on the price tier, % src: results/benchmark_test_task_b.csv [Tree gini (ours), F1 0.9258]
and floor area and location explain most of its decisions. On validation
the linear SVM narrowly beat the Fourier-feature SVM, and only the linear
SVM is stable when trained on less data: the RFF configuration with the
best validation F1 needs more Pegasos updates than a cross-validation
fold provides. The remaining errors sit at
the tier threshold, at the price extremes and in districts with few
listings, and a few come from implausible listings. The price of writing
the tree ourselves is speed: it fits about 277 times more slowly. % src: 86.7520 / 0.3128 from results/benchmark_test_task_b.csv
Future work: ensembles of our trees (bagging or boosting), a vectorised
split search over sorted values with running sums, which would also make
a proper tree sweep affordable, and better location features, such as
distance to the city centre or target statistics of each district
computed on training folds.

\section*{Author Contributions}
{\raggedright
Mahammad Mammadli implemented the data pipeline (\texttt{data\_prep.py}):
loading, cleaning, parsing of floor and land area, leakage removal,
encoding, the stratified split and the training-only tier threshold.
Ibrahim Aghazade implemented the from-scratch decision tree
(\texttt{decision\_tree.py}) with Gini, entropy and MSE criteria, the four
stopping parameters, node statistics and the scikit-learn match test, and
the tree hyperparameter study. Sahib Aliyev implemented the Pegasos SVM
(\texttt{svm.py}) with margin, support vectors and random Fourier
features, plus the sweep and plotting tools (\texttt{svm\_experiment.py},
\texttt{svm\_plots.py}) and their tests. Khizri Sakhrayev implemented the
NumPy metrics and tree feature importances (\texttt{evaluate.py}), the
scikit-learn benchmark (\texttt{benchmark.py}), the SVM tuning on the real
data with the stability-based selection rule (\texttt{svm\_tuning.py}),
the error analysis (\texttt{analysis.py}), the EDA figures
(\texttt{eda.py}), the one-command pipeline (\texttt{run\_all.py}) and the
README, fixed the encoding of the five binary flag columns in the data
pipeline, and wrote and assembled the report.\par}

\section*{Tools and Acknowledgements}
{\raggedright
Khizri Sakhrayev used Claude Code (Anthropic) to produce
\texttt{evaluate.py}, \texttt{tests/test\_evaluate.py},
\texttt{benchmark.py}, \texttt{svm\_tuning.py}, \texttt{analysis.py},
\texttt{eda.py}, \texttt{splits.py}, \texttt{run\_all.py}, the README
and the text of this report, under his direction and review. Sahib Aliyev
used Codex for the SVM implementation, random Fourier features, tuning and
plotting helpers, tests and implementation documentation, as recorded in
\texttt{docs/SVM.md}. We used NumPy, pandas, matplotlib and
scikit-learn~\cite{sklearn}; scikit-learn only for the baselines and
tests.\par}

% =====================================================================
\begin{thebibliography}{9}
\bibitem{cart}
L.~Breiman, J.~H. Friedman, R.~A. Olshen, and C.~J. Stone,
\emph{Classification and Regression Trees}. Belmont, CA: Wadsworth, 1984.

\bibitem{cortes}
C.~Cortes and V.~Vapnik, ``Support-vector networks,''
\emph{Machine Learning}, vol.~20, no.~3, pp.~273--297, 1995.

\bibitem{pegasos}
S.~Shalev-Shwartz, Y.~Singer, N.~Srebro, and A.~Cotter,
``Pegasos: primal estimated sub-gradient solver for SVM,''
\emph{Mathematical Programming}, vol.~127, no.~1, pp.~3--30, 2011.

\bibitem{rahimi}
A.~Rahimi and B.~Recht, ``Random features for large-scale kernel
machines,'' in \emph{Advances in Neural Information Processing Systems
20}, 2007, pp.~1177--1184.

\bibitem{sklearn}
F.~Pedregosa \emph{et al.}, ``Scikit-learn: Machine learning in
Python,'' \emph{Journal of Machine Learning Research}, vol.~12,
pp.~2825--2830, 2011.

\bibitem{kaggle}
sehriyarmemmedli (Kaggle user), ``binaaz-sale-project,'' Kaggle dataset,
2025.
\url{https://www.kaggle.com/datasets/sehriyarmemmedli/binaaz-sale-project}
\end{thebibliography}

% =====================================================================
\onecolumn
\appendix
\renewcommand{\thesection}{\Alph{section}}
\section{Full Results}
\label{sec:appendix}
Tables~\ref{tab:full-a} and~\ref{tab:full-b} give every metric of every
model: the held-out test value and, below it, the 5-fold CV mean $\pm$
std. AZN errors are in thousands of AZN. $^{*}$Fit on at most 10{,}000
training rows.

\begin{table}[h]
  \centering
  \small
  \caption{Task~A (regression of $\log(\text{price})$), all metrics.}
  \label{tab:full-a}
  \begin{tabular}{llrrrrrr}
\toprule
Model & Eval. & RMSE (log) & MAE (log) & $R^2$ (log) & RMSE (kAZN) & MAE (kAZN) & $R^2$ (AZN) \\
\midrule
Tree (ours) & test & 0.241 & 0.166 & 0.853 & 124.9 & 51.5 & 0.753 \\
 & CV & 0.257 $\pm$ 0.006 & 0.172 $\pm$ 0.003 & 0.831 $\pm$ 0.006 & 139.6 $\pm$ 11.1 & 54.3 $\pm$ 1.2 & 0.674 $\pm$ 0.037 \\
Tree (sklearn) & test & 0.242 & 0.166 & 0.852 & 128.4 & 51.7 & 0.739 \\
 & CV & 0.256 $\pm$ 0.005 & 0.172 $\pm$ 0.003 & 0.831 $\pm$ 0.005 & 140.8 $\pm$ 12.1 & 54.3 $\pm$ 1.0 & 0.669 $\pm$ 0.034 \\
\bottomrule
\end{tabular}

\end{table}

\begin{table}[h]
  \centering
  \small
  \caption{Task~B (premium tier), all metrics for the premium class.}
  \label{tab:full-b}
  \begin{tabular}{llrrrr}
\toprule
Model & Eval. & Precision & Recall & F1 & ROC-AUC \\
\midrule
Tree gini (ours) & test & 0.920 & 0.922 & 0.921 & 0.970 \\
 & CV & 0.907 $\pm$ 0.009 & 0.914 $\pm$ 0.006 & 0.911 $\pm$ 0.003 & 0.959 $\pm$ 0.003 \\
Tree gini (sklearn) & test & 0.920 & 0.922 & 0.921 & 0.970 \\
 & CV & 0.907 $\pm$ 0.008 & 0.913 $\pm$ 0.005 & 0.910 $\pm$ 0.003 & 0.958 $\pm$ 0.003 \\
Tree entropy (ours) & test & 0.930 & 0.906 & 0.918 & 0.974 \\
 & CV & 0.903 $\pm$ 0.011 & 0.918 $\pm$ 0.007 & 0.910 $\pm$ 0.003 & 0.963 $\pm$ 0.003 \\
Tree entropy (sklearn) & test & 0.930 & 0.906 & 0.918 & 0.975 \\
 & CV & 0.902 $\pm$ 0.011 & 0.919 $\pm$ 0.006 & 0.910 $\pm$ 0.003 & 0.964 $\pm$ 0.003 \\
Pegasos linear (ours) & test & 0.900 & 0.898 & 0.899 & 0.960 \\
 & CV & 0.898 $\pm$ 0.008 & 0.902 $\pm$ 0.007 & 0.900 $\pm$ 0.003 & 0.959 $\pm$ 0.003 \\
SGD hinge linear (sklearn) & test & 0.903 & 0.894 & 0.898 & 0.961 \\
 & CV & 0.902 $\pm$ 0.004 & 0.891 $\pm$ 0.009 & 0.896 $\pm$ 0.004 & 0.959 $\pm$ 0.004 \\
SVC linear (sklearn)$^{*}$ & test & 0.896 & 0.898 & 0.897 & 0.959 \\
 & CV & 0.900 $\pm$ 0.004 & 0.905 $\pm$ 0.006 & 0.902 $\pm$ 0.003 & 0.962 $\pm$ 0.003 \\
Pegasos RFF (ours) & test & 0.872 & 0.920 & 0.895 & 0.961 \\
 & CV & 0.887 $\pm$ 0.055 & 0.877 $\pm$ 0.059 & 0.879 $\pm$ 0.010 & 0.958 $\pm$ 0.003 \\
SGD hinge RFF (sklearn) & test & 0.910 & 0.874 & 0.892 & 0.959 \\
 & CV & 0.893 $\pm$ 0.035 & 0.881 $\pm$ 0.044 & 0.885 $\pm$ 0.006 & 0.958 $\pm$ 0.003 \\
SVC RBF (sklearn)$^{*}$ & test & 0.902 & 0.895 & 0.898 & 0.963 \\
 & CV & 0.907 $\pm$ 0.006 & 0.904 $\pm$ 0.006 & 0.905 $\pm$ 0.005 & 0.968 $\pm$ 0.003 \\
\bottomrule
\end{tabular}

\end{table}

\end{document}
